\scanpage{246}
\mspage{196}

qu'à isos \inserted{(peut-être} non canoniques) près, le foncteur composé défini par $C$

\[
X\longmapsto C(X)\longmapsto (s\longmapsto\unsure{\mathrm{hot}|\mathrm{Hot}}_W(C(X)_s))
\]

\diagram{\begin{tikzcd}[ampersand replacement=\&,column sep=1.7em]
\mathrm{Cat}_{/S}\arrow[r]\arrow[rr,bend right=18,"{\unsure{\varphi|\phi}_C}"']\&\underline{\mathrm{Fib}}_W S\arrow[r]\&\underline{\mathrm{Hom}}(S^{\unsure{\circ|o}}\unsure{,|;}\mathrm{Hot}_W)
\end{tikzcd}}

ne dépend pas du choix particulier de $C$, i.e. deux tels foncteurs ``fibres homotopiques'' (correspondant à deux choix $C,C'$) \unsure{sont} toujours isomorphes.

Mais au lieu d'utiliser $C\circ C'$ pour définir un tel isomorphisme, on aurait pu utiliser $C'\circ C$, ou quelque autre procédé. On aimerait comprendre s'il y a unicité véritable, une caractérisation intrinsèque de ces isomorphismes $\varphi_C\xrightarrow{\sim}\varphi_{C'}$, indépendamment de telles constructions ou de telles autres.

Je présume que $C\circ C'$ est tout aussi muni de $\mathrm{id}_{\mathrm{Cat}_{/S}}\longrightarrow C\circ C'$,
% archive sheet number: 246
